%% start of file `template.tex'.
%% Copyright 2006-2015 Xavier Danaux (xdanaux@gmail.com).
%
% This work may be distributed and/or modified under the
% conditions of the LaTeX Project Public License version 1.3c,
% available at http://www.latex-project.org/lppl/.


\documentclass[11pt,a4paper,sans]{moderncv}        % possible options include font size ('10pt', '11pt' and '12pt'), paper size ('a4paper', 'letterpaper', 'a5paper', 'legalpaper', 'executivepaper' and 'landscape') and font family ('sans' and 'roman')

% moderncv themes
\moderncvstyle{casual}                             % style options are 'casual' (default), 'classic', 'banking', 'oldstyle' and 'fancy'
\moderncvcolor{blue}                               % color options 'black', 'blue' (default), 'burgundy', 'green', 'grey', 'orange', 'purple' and 'red'
%\renewcommand{\familydefault}{\sfdefault}         % to set the default font; use '\sfdefault' for the default sans serif font, '\rmdefault' for the default roman one, or any tex font name
%\nopagenumbers{}                                  % uncomment to suppress automatic page numbering for CVs longer than one page

% character encoding
\usepackage[utf8]{inputenc}                       % if you are not using xelatex ou lualatex, replace by the encoding you are using
%\usepackage{CJKutf8}                              % if you need to use CJK to typeset your resume in Chinese, Japanese or Korean

% adjust the page margins
\usepackage[scale=0.75]{geometry}
%\setlength{\hintscolumnwidth}{3cm}                % if you want to change the width of the column with the dates
%\setlength{\makecvtitlenamewidth}{10cm}           % for the 'classic' style, if you want to force the width allocated to your name and avoid line breaks. be careful though, the length is normally calculated to avoid any overlap with your personal info; use this at your own typographical risks...

% personal data
\name{Kundai}{Sachikonye}
%\title{Resumé title}                               % optional, remove / comment the line if not wanted
%\address{Jagdstrasse 13}{85356}% optional, remove / comment the line if not wanted; the "postcode city" and "country" arguments can be omitted or provided empty
%\phone[mobile]{+1~(234)~567~890}                   % optional, remove / comment the line if not wanted; the optional "type" of the phone can be "mobile" (default), "fixed" or "fax"
%\phone[fixed]{+2~(345)~678~901}
%\phone[fax]{+3~(456)~789~012}
\email{kundai.sachikonye@wzw.tum.de}                               % optional, remove / comment the line if not wanted
%\homepage{www.johndoe.com}                         % optional, remove / comment the line if not wanted
%\social[linkedin]{john.doe}                        % optional, remove / comment the line if not wanted
%\social[twitter]{jdoe}                             % optional, remove / comment the line if not wanted
\social[github]{fullscreen-triangle}                              % optional, remove / comment the line if not wanted
%\extrainfo{additional information}                 % optional, remove / comment the line if not wanted
%\photo[64pt][0.4pt]{picture}                       % optional, remove / comment the line if not wanted; '64pt' is the height the picture must be resized to, 0.4pt is the thickness of the frame around it (put it to 0pt for no frame) and 'picture' is the name of the picture file
\quote{Some quote}                                 % optional, remove / comment the line if not wanted

                % 'publications' is the name of a BibTeX file

\clearpage
%-----       letter       ---------------------------------------------------------
% recipient data
\recipient{Universitätsklinikum Münster}{Albert-Schweitzer-Campus 1\\  48149 Münster \\  Germany}
\date{September 19 2024}
\opening{Dear Dr. Sir/Madam}
\closing{Yours faithfully,}
\enclosure[Attached]{curriculum vit\ae{}}          % use an optional argument to use a string other than "Enclosure", or redefine \enclname
\makelettertitle

I am writing to express my enthusiasm for a job in bioinformatics and data management. With a diverse background spanning biochemistry, pharmaceutical biotechnology, and computational systems medicine, I bring a unique blend of skills and experiences to the table. 

In order for me to achieve the goals of my studies in pharmaceutical biotechnology, I had to familiarise myself with event driven software architecture. The automation of bioprocesses required the use of distributed  hybrid service oriented server architecture in multiple languages that leveraged web frameworks which exposed remote,web, transaction APIs, incoming or outgoing HTTP/TCP and socket connections.
All the downstream modules involved in signal processing, data analysis and process control were integrated as loosely coupled headless server modules which left ample room for further detailed implementations isolated from the main deployment. The growth kinetics could be modeled as both ordinary or stochastic processes, allowing application of reinforcement learning, dynamic and linear programming methods. Systems biology methods were used to characterise pathways that were subsequently used in conjunction with denovo microbiome and proteomic assembly algorithms in predicting microbial communities.

In the realm of computational lipidomics, my task was similar in its goals albeit with differences in the scope of the pipeline. This meant the development of a high throughput cluster-based distributed hybrid backend server solution with an all inclusive policy which accepted all mass spectrometry based lipidomics data, regardless of vendor, instrument configuration or acquisition type and executed standard data processing pipelines. The global nature of the project introduced restrictions on data access, task scheduling, cross-compatability and the abstraction of database queries through the use of multiple  authorization,  tokenization, serialization within the FAIR foundational principals of research data ecosystems. Knowledge of bioinformatics  was needed for complete comprehension of molecular representations, the application of algorithms for common substructure searches, structure-activity relationships, fingerprinting and fragmentation trees. This extended from gene based ontologies in whole genome denovo assembly, computational models in declarative systems biology markup languages in drug re-purposing
and engineering of fragment ions into features for deep learning.


At Bitspark, my proficiency was extended beyond the domain of academic research and I was transformed into an individual capable of developing complete annotation frameworks for large datasets from virtually an field.I have managed to develop serverless, headless and stateless modules for communication protocols with GNSS chipsets and smart device sensors via Bluetooth, WIFI or GMS orchestrated by asynchronous functions forming a cross-device mesh network of IoT devices.  The privacy, resource and distribution constraints of smart devices are perfectly fulfilled by federated learning which allows for local data training whilst transmitting only updated model parameters in a process that is highly scalable.


I have always oriented my life around a secure orientation towards novelty,exploration, growth, things I am eager to discover at your workplace. Thank you for considering my application. 







  
\makeletterclosing

%\clearpage\end{CJK*}                              % if you are typesetting your resume in Chinese using CJK; the \clearpage is required for fancyhdr to work correctly with CJK, though it kills the page numbering by making \lastpage undefined
\end{document}


%% end of file `template.tex'.
